\section{Diseño general del estudio}
\label{sec:disenio_general}

El modelo predictivo se construyó siguiendo un enfoque de aprendizaje automático basado en una arquitectura LSTM multi-torre, la cual procesa series temporales de predictores climáticos de múltiples variables y resoluciones espaciales para estimar la precipitación mensual acumulada (en mm) en una región de interés. El flujo de trabajo se estructura en cuatro etapas principales: (1) adquisición y preprocesamiento de datos, (2) normalización y reducción de dimensionalidad, (3) entrenamiento y validación del modelo, y (4) evaluación mediante diagnósticos completos.

\section{Adquisición y preprocesamiento de datos}
\label{sec:adquisicion_datos}

\subsection{Fuentes de datos}\label{subsec:fuentes_de_datos}
Los datos de precipitación observada fueron obtenidos del Instituto Nacional de Sismología, Vulcanología, Meteorología e Hidrología (INSIVUMEH) de Guatemala, consistiendo en series de tiempo diarias de precipitación medida en estaciones pluviométricas calibradas distribuidas a lo largo del país. Los predictores climáticos fueron adquiridos de fuentes globales de reanálisis (Copernicus Climate Data Store: ERA5 post-processed daily statistics on single [and pressure] levels from 1940 to present) considerados en la cuadrícula de latitudes de $[-165, 0]$ en longitud y $[-15, 75]$ en latitud, para poder considerar las teleconexiones que dominan el escenario climático guatemalteco, sin considerar ruido proveniente de otras regiones climáticas lejos de Centroamérica:

\begin{itemize}
    \item Temperatura Superficial del Mar (SST).
    \item Presión Superficial (SP).
    \item Humedad relativa y específica (r y q) a 500 hPa. 
    \item Geopotencial (z) a 500 hPa. 
    \item Velocidades zonal y meridional (u, v) a 850 hPa. 
    \item Índices de Madden-Julian Oscillation (MJO) (RMM1, RMM2).
    \item Precipitación global (tp).
    \item Temperatura a 2 metros de la superficie (2m).
\end{itemize}

Se asume que determinadas combinaciones de predictores tendrán mejores resultados que otras, tanto para el país en general, como para regiones y departamentos específicos de este. Las combinaciones propuestas fueron:
\begin{enumerate}
    \item SST, z, r, q, tp, u, v.
    \item SST, SP, tp.
    \item SST, SP, tp, MJO.
    \item SST, tp, MJO. 
    \item SST, SP, z, r, q, tp, u, v.
    \item SST, SP, z, r, q, tp, u, v, MJO.
    \item SST, tp.
    \item SST, sp, tp, 2t
    \item SST, tp, 2t
    \item SST, tp, MJO, 2t
    \item SST, z, r, q, tp, u, v, 2t.
    \item SST, z, r, q, tp, u, v, 2t, MJO.
\end{enumerate}

La efectividad de cada predictor será evaluada mediante diagnósticos descritos más adelante.

\subsection{Normalización espacio-temporal}
Todos los archivos de entrada (NetCDF4) fueron procesados para garantizar consistencia dimensional y nomenclatura:

\begin{itemize}
    \item \textbf{Dimensiones:} Se normalizó la nomenclatura de ejes a \texttt{(time, lat, lon)}, independientemente del esquema original del archivo (\texttt{T}, \texttt{valid\_time}, \texttt{time\_counter}, etc.).
    \item \textbf{Coordenadas:} Se identificaron automáticamente los nombres de latitud y longitud (\texttt{latitude}, \texttt{y}, \texttt{Y}) y se renombraron a \texttt{lat} y \texttt{lon}.
    \item \textbf{Indexación temporal:} Se convirtieron valores numéricos o cadenas de tiempo a formato \texttt{datetime64} de NumPy para facilitar operaciones de fecha y sincronización.
\end{itemize}

\subsection{Detección de píxeles válidos y máscaras espaciales}
Dado que las grillas rasterizadas típicamente contienen píxeles con NaN (océanos, regiones fuera de dominio), se construyó una máscara de validez mediante muestreo de los primeros 50 pasos temporales de cada predictor. Un píxel se clasificó como válido si presentaba al menos un valor finito en el intervalo muestreado. Esta máscara aseguró que:

\begin{itemize}
    \item La red neuronal procesa solo píxeles con datos físicamente significativos.
    \item Se evita acumular gradientes en píxeles sin información.
    \item La compresión dimensional mediante PCA se realiza exclusivamente sobre píxeles válidos.
\end{itemize}

\section{Normalización y reducción de dimensionalidad}
\label{sec:normalizacion_pca}

\subsection{Estandarización de predictores}
Para que el optimizador por gradiente converja eficientemente, cada predictor fue estandarizado usando estadísticas (media $\mu_p$ y desviación estándar $\sigma_p$) calculadas sobre el conjunto de entrenamiento:

\begin{equation}
    X_{\text{std}}^{(p)} = \frac{X^{(p)} - \mu_p}{\sigma_p}
\end{equation}

donde $p$ indexa cada uno de los $n_{\text{pred}}$ predictores. Estas estadísticas se almacenaron en la metadata del entrenamiento para reproducibilidad y aplicación consistente en predicciones posteriores.

\subsection{Análisis de Componentes Principales (PCA)}
Como los predictores climáticos típicamente presentan elevada colinealidad espacial (continuidad suave de campos como temperatura y presión), se aplicó PCA para reducir la dimensionalidad:

\begin{enumerate}
    \item Se calculó la matriz de covarianza de cada predictor estandarizado sobre el conjunto de entrenamiento en píxeles válidos.
    \item Se seleccionaron las primeras $k_p$ componentes principales que acumulaban al menos el 99.9\% de la varianza explicada, siempre que fuera posible.
    \item El modelo PCA se almacenó para transformación consistente durante predicción.
\end{enumerate}

La reducción dimensional cumple varios objetivos:
\begin{itemize}
    \item Disminuye el número de parámetros de la red neuronal, mitigando el sobreajuste.
    \item Comprime información correlacionada en componentes ortogonales.
    \item Acelera el cálculo durante entrenamiento y predicción.
\end{itemize}

\section{Construcción y segmentación de secuencias temporales}
\label{sec:secuencias_temporales}

\subsection{Ventanas de entrada}
Para la predicción de precipitación mensual, la red LSTM requiere una ventana de entrada de longitud fija $L$ (por defecto, $L = 31$ días). Formalmente, si se predice la precipitación acumulada en el mes $[t_{\text{start}}, t_{\text{end}}]$, la secuencia de entrada comprende $L$ días anteriores:

\begin{equation}
    \mathbf{x}_{t-L:t} = [X_{t-L}, X_{t-L+1}, \ldots, X_{t-1}]
\end{equation}

donde cada $X_{\tau}$ es un vector de características (predictores normalizados y comprimidos por PCA) para el día $\tau$.

\subsection{División cronológica del conjunto de datos}
Para evaluar la capacidad de generalización del modelo sin contaminación temporal (data leakage), el conjunto de datos se dividió de manera cronológica en dos subconjuntos:

\begin{itemize}
    \item \textbf{Entrenamiento:} Años iniciales (80\% del período total). Utilizado para ajustar los pesos $\mathbf{W}$ y sesgos $\mathbf{b}$ de la red mediante retropropagación del error.
    \item \textbf{Validación:} Período intermedio (20\%). Monitorea el desempeño fuera de muestra durante el entrenamiento, permitiendo la detención temprana (\textit{early stopping}) para evitar sobreajuste.
\end{itemize}

Esta partición respeta la estructura temporal inherente a los datos climáticos, donde la dependencia estadística es más fuerte entre observaciones cercanas en tiempo.

\section{Arquitectura del modelo LSTM multi-torre}
\label{sec:arquitectura_lstm}

\subsection{Diseño de entrada y torres paralelas}
El modelo multi-torre procesa cada predictor climático mediante una rama (torre) independiente de capas LSTM y capas totalmente conectadas, explotando la especificidad de cada variable:

\begin{enumerate}
    \item \textbf{Torre $p$:} Recibe una secuencia de entrada $\mathbf{x}_{t-L:t}^{(p)} \in \mathbb{R}^{L \times d_p}$, donde $d_p$ es la dimensión post-PCA del predictor $p$.
    \item \textbf{Capas LSTM:} La secuencia temporal se procesa mediante una o más capas LSTM que capturan dependencias de largo alcance. Cada LSTM emite un estado oculto $\mathbf{h}_T^{(p)}$ al final de la ventana.
    \item \textbf{Capas densas locales:} El estado final se procesa mediante capas totalmente conectadas para extraer características no lineales específicas del predictor.
\end{enumerate}

\subsection{Fusión de representaciones y decodificación espacial}
Las salidas de todas las torres se concatenan:

\begin{equation}
    \mathbf{z} = [\mathbf{z}^{(1)}, \mathbf{z}^{(2)}, \ldots, \mathbf{z}^{(n_{\text{pred}})}]
\end{equation}

A continuación, se aplican varias técnicas de refinamiento espacial:

\begin{enumerate}
    \item \textbf{Modulación por coordenadas:} Un módulo aprendible que ingiere las coordenadas geográficas (lat, lon) de cada píxel y genera factores de escala y desplazamiento (shift) para modular las predicciones. Esta técnica captura variabilidad geográfica implícita.
    
    \item \textbf{Sesgo por píxel:} Se añade un sesgo aprendible específico de cada píxel, permitiendo que el modelo capture errores sistemáticos locales.
    
    \item \textbf{Refinamiento espacial:} Un módulo convolucional (2D) que reconstruye temporalmente la grilla $(h_{\text{target}} \times w_{\text{target}})$ a partir de píxeles válidos, propaga información espacial mediante capas convolucionales (Conv2D + BatchNorm), y reextrae los píxeles tierra para corrección fina.
\end{enumerate}

\subsection{Transformación de objetivo}
Debido a la distribución fuertemente asimétrica (sesgada) de la precipitación mensual, se aplicó una transformación invertible al objetivo antes del entrenamiento. Las opciones implementadas incluyen:

\begin{itemize}
    \item \textbf{Log-transformación:} $y_{\text{trans}} = \log(y + 1)$, adecuada para datos con rango amplio.
    \item \textbf{Raíz cúbica:} $y_{\text{trans}} = y^{1/3}$, reduce el impacto de valores extremos manteniendo el orden.
    \item \textbf{Raíz cuarta:} $y_{\text{trans}} = y^{1/4}$, estabilización más fuerte.
    \item \textbf{Raíz cuadrada:} $y_{\text{trans}} = \sqrt{y}$, apropiada para datos con varianza proporcional a la media.
    \item \textbf{Estandarización:} $y_{\text{trans}} = \frac{y - \mu_y}{\sigma_y}$, para varianza homogénea.
\end{itemize}

El modelo LSTM entrenó en el espacio transformado, y las predicciones se invirtieron analíticamente para recuperar precipitación en milímetros.

\section{Entrenamiento del modelo}
\label{sec:entrenamiento_modelo}


\subsection{Función de Pérdida Compuesta y Argumentos de Entrenamiento}

El modelo utiliza una función de pérdida compuesta que integra múltiples componentes diseñados para capturar diferentes aspectos del problema de predicción de precipitación. A continuación se describen las funciones de pérdida implementadas y los parámetros principales utilizados durante el entrenamiento.

\subsubsection{Componentes de la Función de Pérdida}

La pérdida total es una combinación ponderada de cuatro componentes principales:

\begin{equation}
	\mathcal{L}_{\text{total}} = w_H \cdot \mathcal{L}_{\text{Huber}} + w_Q \cdot \mathcal{L}_{\text{Quantil}} + w_T \cdot \mathcal{L}_{\text{Tweedie}} + w_C \cdot \mathcal{L}_{\text{Correlación}} + \mathcal{L}_{\text{Suavizado}}
\end{equation}

donde $w_H$, $w_Q$, $w_T$ y $w_C$ son los pesos relativos de cada componente, relacionados mediante:
\begin{equation}
	w_H = 1.0 - w_Q - w_T - w_C
\end{equation}


\paragraph{1. Pérdida de Huber Focal} 
La función de Huber es robusta ante valores atípicos. Se define como:
\begin{equation}
	\mathcal{L}_{\text{Huber}}(y, \hat{y}) = \begin{cases}
		\frac{1}{2}(y - \hat{y})^2 & \text{si } |y - \hat{y}| \leq \delta \\
		\delta \left(|y - \hat{y}| - \frac{\delta}{2}\right) & \text{si } |y - \hat{y}| > \delta
	\end{cases}
\end{equation}

El parámetro delta (\texttt{--huber\_delta}) controla el punto de transición entre comportamiento cuadrático y lineal. La versión focal amplifica el error en eventos extremos mediante un factor $w_{\text{focal}}$ que depende de la intensidad de precipitación:

\begin{equation}
	w_{\text{focal}} = 1.0 + (\text{extreme\_boost} - 1.0) \cdot \sigma\left(\frac{y_{\text{mm}} - h_{\text{high}}}{h_{\text{width}}}\right)
\end{equation}

donde $\sigma$ es la función sigmoide, $h_{\text{high}}$ es el umbral superior en mm, y \texttt{--extreme\_boost} controla la amplificación.

\paragraph{2. Pérdida de Cuantil Adaptativa}
La pérdida de cuantil se adapta según la intensidad de precipitación, utilizando tres regímenes:

\begin{equation}
	\mathcal{L}_{\text{Quantil}} = \mathbb{E}\left[w_{\text{low}} \cdot \ell_{\tau=0.50} + w_{\text{mid}} \cdot \ell_{\tau=0.60} + w_{\text{high}} \cdot \ell_{\tau_{\text{user}}}\right]
\end{equation}

donde $\ell_\tau(y, \hat{y}) = \max(\tau(y - \hat{y}), (\tau - 1)(y - \hat{y}))$ es la pérdida de cuantil estándar.

Las regiones se definen mediante transiciones suaves (sigmoide):
\begin{itemize}
	\item \textbf{Zona baja} ($y < \texttt{--low\_threshold}$ mm): $\tau = 0.50$ (mediana), captura eventos débiles
	\item \textbf{Zona media} ($\texttt{--low\_threshold} \leq y \leq \texttt{--high\_threshold}$): $\tau = 0.60$, transición intermedia
	\item \textbf{Zona alta} ($y > \texttt{--high\_threshold}$ mm): $\tau = \texttt{--quantile\_tau}$, captura eventos extremos
\end{itemize}

Los pesos $w_{\text{low}}$, $w_{\text{mid}}$, $w_{\text{high}}$ se calculan mediante normalización softmax de las salidas sigmoides para garantizar $\sum w_i = 1$.

\paragraph{3. Pérdida de Tweedie}
La distribución Tweedie es una distribución compuesta Poisson-Gamma natural para datos de precipitación. La pérdida es:

\begin{equation}
	\mathcal{L}_{\text{Tweedie}} = \mathbb{E}\left[-y_{\text{mm}} \cdot \hat{y}_{\text{mm}}^{-(1-p)} / (1-p) + \hat{y}_{\text{mm}}^{(2-p)} / (2-p)\right] / s_{\text{tw}}
\end{equation}

donde $p = \texttt{--tweedie\_power} \in (1, 2)$ interpola entre Poisson ($p=1$) y Gamma ($p=2$). El parámetro $s_{\text{tw}} = 1.0 / (2 - p)$ normaliza la escala cuando \texttt{--tweedie\_scale\_norm} es verdadero, mejorando la estabilidad numérica.

\paragraph{4. Pérdida de Correlación de Pearson}
Maximiza la correlación espacial entre predicciones y observaciones:

\begin{equation}
	\mathcal{L}_{\text{Correlación}} = 1.0 - r_{\text{Pearson}}
\end{equation}

donde:
\begin{equation}
	r_{\text{Pearson}} = \frac{\text{Cov}(\hat{y}, y)}{\sigma_{\hat{y}} \cdot \sigma_y}
\end{equation}

\paragraph{5. Pérdida de Suavizado Espacial}
Penaliza diferencias grandes entre predicciones de píxeles adyacentes, promoviendo coherencia espacial:

\begin{equation}
	\mathcal{L}_{\text{Suavizado}} = \texttt{--smooth\_weight} \cdot \mathbb{E}\left[\sum_{\text{pares vecinos}} (\hat{y}_i - \hat{y}_j)^2\right]
\end{equation}

\subsection{Regularización y detención temprana}
Para mitigar sobreajuste en esta red profunda se aplicaron:

\begin{enumerate}
    \item \textbf{Dropout:} Se desactivó aleatoriamente un porcentaje (típicamente 20-30\%) de neuronas en cada paso durante entrenamiento, forzando redundancia y robustez.
    \item \textbf{Batch Normalization:} Se normalizaron las activaciones por minilote, acelerando convergencia y mejorando generalización.
    \item \textbf{Early Stopping:} El entrenamiento se detuvo cuando la métrica de validación (RMSE o KGE) no mejoró durante $N$ épocas consecutivas (típicamente $N=15$), previniendo descenso en generalización.
\end{enumerate}

\subsection{Acumulación de gradientes}
Para facilitar el entrenamiento con lotes grandes (que mejora estabilidad del gradiente), se implementó acumulación de gradientes: cada paso acumula gradientes de múltiples minilotes (típicamente 4-8) antes de actualizar pesos. Esto efectivamente incrementa el tamaño del lote sin aumentar el uso de memoria GPU.

\subsubsection{Normalización EMA de componentes}

Durante el entrenamiento, las magnitudes relativas de cada componente pueden variar significativamente. Para mantener equilibrio automático, el modelo utiliza medias móviles exponenciales (EMA) que se actualizan solo durante los primeros 50 pasos de entrenamiento (fase de \textit{warmup}):

\begin{equation}
\text{EMA}_i^{(t)} = \begin{cases}
\mathcal{L}_i & \text{si } t = 0 \\
\alpha \cdot \text{EMA}_i^{(t-1)} + (1 - \alpha) \cdot \mathcal{L}_i & \text{si } 0 < t \leq 50 \\
\text{EMA}_i^{(50)} & \text{si } t > 50
\end{cases}
\end{equation}

donde $\alpha = 0.95$ es el momentum. Cada componente se normaliza dividiéndola por su EMA estática:

\begin{equation}
\mathcal{L}_{\text{total}} = w_H \cdot \frac{\mathcal{L}_{\text{Huber}}}{\text{EMA}_H} + w_Q \cdot \frac{\mathcal{L}_{\text{Quantil}}}{\text{EMA}_Q} + \cdots
\end{equation}

Esto garantiza que componentes de escalas muy diferentes contribuyan equitativamente a los gradientes.

\subsubsection{Argumentos principales de entrenamiento}

A continuación se describen los argumentos de línea de comandos más relevantes que controlan el entrenamiento del modelo:

\subsection*{Pérdida -- Componente Huber}

\begin{itemize}
    \item \texttt{--huber\_delta}: Punto de transición del criterio de Huber entre pérdida cuadrática y lineal (en unidades transformadas). Valor por defecto: $1.0$.
    
    \item \texttt{--extreme\_boost}: Factor multiplicativo para amplificar gradientes en eventos extremos. Valor por defecto: $3.5$.
\end{itemize}

\subsection*{Pérdida -- Componente Quantil}

\begin{itemize}
    \item \texttt{--quantile\_weight}: Peso relativo de la pérdida de cuantil en la pérdida total. Valor por defecto: $0.30$.
    
    \item \texttt{--quantile\_tau}: Valor de cuantil $\tau$ usado en la zona de intensidad alta. Valor por defecto: $0.72$.
    
    \item \texttt{--low\_threshold}: Umbral inferior de precipitación en mm (frontera zona baja/media). Valor por defecto: $50.0$.
    
    \item \texttt{--high\_threshold}: Umbral superior de precipitación en mm (frontera zona media/alta). Valor por defecto: $250.0$.
\end{itemize}

\subsection*{Pérdida -- Componente Tweedie}

\begin{itemize}
    \item \texttt{--tweedie\_weight}: Peso relativo de la pérdida de Tweedie. Valor por defecto: $0.40$.
    
    \item \texttt{--tweedie\_power}: Parámetro $p \in (1,2)$ de la distribución Tweedie. Valor por defecto: $1.30$.
    
    \item \texttt{--tweedie\_scale\_norm}: Si verdadero, normaliza la escala de Tweedie por el factor $1/(2-p)$. Valor por defecto: \texttt{true}.
\end{itemize}

\subsection*{Pérdida -- Componentes Adicionales}

\begin{itemize}
    \item \texttt{--corr\_weight}: Peso relativo de la pérdida de correlación de Pearson. Valor por defecto: $0.10$.
    
    \item \texttt{--smooth\_weight}: Peso relativo de la pérdida de suavizado espacial. Valor por defecto: $0.005$.
\end{itemize}

\subsection*{Arquitectura y Entrenamiento}

\begin{itemize}
    \item \texttt{--spatial\_embed}: Dimensionalidad inicial de proyección espacial en cada torre. Valor por defecto: $256$.
    
    \item \texttt{--temporal\_dim}: Dimensionalidad de las capas BiLSTM. Valor por defecto: $128$.
    
    \item \texttt{--hidden\_dim}: Dimensionalidad de las capas Dense del decodificador. Valor por defecto: $256$.
    
    \item \texttt{--n\_attn\_heads}: Número de cabezas en el mecanismo de multi-head attention. Valor por defecto: $8$.
    
    \item \texttt{--pixel\_embed\_dim}: Dimensionalidad del embedding por píxel para ajuste fino. Valor por defecto: $64$.
    
    \item \texttt{--coord\_embed\_dim}: Dimensionalidad de modulación por coordenadas geográficas. Valor por defecto: $32$.
\end{itemize}

\subsection*{Datos y Preprocesamiento}

\begin{itemize}
    \item \texttt{--input\_days}: Número de días históricos como entrada temporal. Valor por defecto: $31$.
    
    \item \texttt{--batch\_size}: Tamaño de batch para entrenamiento (ajustado automáticamente si excede RAM). Valor por defecto: $64$.
    
    \item \texttt{--epochs}: Número máximo de épocas de entrenamiento. Valor por defecto: $500$.
    
    \item \texttt{--target\_transform}: Transformación del target: \texttt{log1p}, \texttt{cbrt}, \texttt{pow025}, \texttt{pow05}, \texttt{standard}, \texttt{none}. Valor por defecto: \texttt{log1p}.
    
    \item \texttt{--pca\_components}: Componentes PCA: $>0$ = fijo, $0$ = automático (retiene varianza), $-1$ = sin PCA. Valor por defecto: $1500$.
    
    \item \texttt{--pca\_variance}: Varianza explicada objetivo para PCA automático. Valor por defecto: $0.999$.
\end{itemize}

\subsection*{Regularización y aumento de datos \textit{(data augmentation})}

\begin{itemize}
    \item \texttt{--noise\_std}: Desviación estándar del ruido gaussiano aditivo en augmentation. Valor por defecto: $0.05$.
    
    \item \texttt{--mixup\_alpha}: Parámetro $\alpha$ de la distribución Beta para mezcla de datos. Valor por defecto: $0.40$.
    
    \item \texttt{--l2\_output}: Regularización L2 en la capa de salida. Valor por defecto: $10^{-5}$.
    
    \item \texttt{--grad\_accum\_steps}: Pasos de acumulación de gradientes (simula batch mayor). Valor por defecto: $1$.
\end{itemize}

\subsection*{Optimización y aprendizaje}

\begin{itemize}
    \item \texttt{--use\_swa}: Activar Stochastic Weight Averaging en la última fase de entrenamiento. Valor por defecto: \texttt{false}.
    
    \item \texttt{--prefetch}: Número de batches a precargar \texttt{tf.data}. Valor por defecto: $2$.
    
    \item \texttt{--max\_ram\_gb}: Límite de RAM para el proceso (control de bajo nivel). Valor por defecto: $4.0$.
\end{itemize}

\subsection*{Reproducibilidad}

\begin{itemize}
    \item \texttt{--seed}: Semilla para reproducibilidad global (Python, NumPy y TensorFlow). Valor por defecto: $42$.
    
    \item \texttt{--seed\_split}: Semilla independiente para la división train/validation. Valor por defecto: igual a \texttt{--seed}.
\end{itemize}
\subsubsection{Interacción de parámetros y recomendaciones}

\begin{itemize}
    
    \item \textbf{Balance de la pérdida}: Si se ajustan $w_Q$, $w_T$ o $w_C$, la suma $w_Q + w_T + w_C$ debe ser estrictamente menor a 1.0 para garantizar $w_H > 0$. El código genera un error si esta condición no se cumple.
    
    \item \textbf{Umbrales de quantil}: Los valores de \texttt{--low\_threshold} y \texttt{--high\_threshold} deben ser representativos del rango de precipitación esperado. Valores típicos están en el rango 30-150 mm (baja), 100-300 mm (alta). Las transiciones suaves (sigmoide) garantizan cambios continuos.
    
    \item \textbf{Potencia de distribución Tweedie}: El valor $p \in (1, 2)$ controla el sesgo entre Poisson ($p \approx 1$, enfatiza ceros) y Gamma ($p \approx 2$, enfatiza cola).
    
    \item \textbf{Auto-batch}: Si \texttt{--batch\_size} excede la memoria disponible calculada por el monitor RAM, se reduce automáticamente y se incrementa \texttt{--grad\_accum\_steps} para mantener un batch efectivo similar. 
    
    \item \textbf{Spatial embed $\gg$ PCA}: Si $w_Q + w_T + w_C < 0.5$, la pérdida de Huber domina. En ese caso, aumentar \texttt{--huber\_delta} (más tolerancia a valores atípicos) o \texttt{--extreme\_boost} (más énfasis en valores extremos).
\end{itemize}

\subsubsection{Implementación de Control de Memoria (Low-RAM)}

El modelo implementa un monitoreo dinámico de RAM en tiempo real mediante \texttt{RAMMonitor}:

\begin{itemize}
    \item \textbf{Límite configurable}: El argumento \texttt{--max\_ram\_gb} especifica la RAM máxima que el proceso puede usar. El monitor mantiene un margen de seguridad del 20\%.
    
    \item \textbf{Lectura desde disco}: Los predictores y target nunca se cargan completamente en la memoria RAM. Se utilizan memmap de NumPy en modo lectura que aprovechan el \textit{page cache} del sistema operativo.
    
    \item \textbf{Procesamiento por chunks}: Las operaciones de estadística, PCA, y predicción se procesan en chunks temporales cuyo tamaño se ajusta dinámicamente según la RAM disponible.
    
    \item \textbf{Liberado de memoria}: Se ejecuta \texttt{gc.collect()} después de cada etapa intensiva para liberar memoria inmediatamente.
\end{itemize}

Esta arquitectura permite entrenar modelos con predictores de cientos de miles de píxeles incluso en máquinas con poca RAM, repartiendo la carga entre procesamiento en disco (I/O rápido) y cálculo (GPU/CPU).


\section{Evaluación del modelo}
\label{sec:evaluacion_modelo}

\subsection{Métricas de error global}
El desempeño predictivo se cuantificó mediante:

\begin{enumerate}
    \item RMSE (Raíz del error cuadrático medio).
    \item MAE (Error Absoluto Medio).
    \item Coeficiente de Determinación ($R^2$).
    \item Correlación de Pearson ($r$).
\end{enumerate}

\subsection{Métricas de eficiencia de modelos hidrológicos}
Se incluyeron las métricas estándar en la predicción de precipitación:

\begin{enumerate}
    \item Nash-Sutcliffe Efficiency (NSE).
    \item Kling-Gupta Efficiency (KGE). 
    \item Porcentaje de Sesgo (PBIAS):
    \begin{equation}
        \text{PBIAS} = 100 \cdot \frac{\sum (\hat{y}_i - y_i)}{\sum y_i}
    \end{equation}
    Cuantifica si el modelo tiende a subestimar o sobreestimar la precipitación.
\end{enumerate}

\subsection{Análisis temporal y espacial}
Para diagnosticar fortalezas y debilidades del modelo:

\begin{enumerate}
    \item \textbf{Variación mensual:} Se calcularon RMSE, bias, correlación y KGE para cada mes del año, identificando estaciones de mejor/peor desempeño.
    \item \textbf{Variación interanual:} Métricas agregadas por año hidrológico permiten detectar épocas de reducida habilidad predictiva (ej. bajo impacto de ENSO).
    \item \textbf{Análisis de extremos:} Se evaluó desempeño específicamente en percentiles altos (P90, P95, P99) de precipitación, críticos para alerta temprana.
    \item \textbf{Mapas espaciales:} Se generaron grillas de bias, RMSE, correlación, NSE y KGE píxel a píxel, revelando heterogeneidad geográfica del error.
\end{enumerate}

\subsection{Distribuciones y estadísticas de entrenamiento}

Con el propósito de complementar, se consideró lo siguiente:
\begin{enumerate}
    \item \textbf{QQ-plot:} Gráfico de cuantiles observados vs predichos, que expone desviaciones sistemáticas en colas.
    \item \textbf{Función de distribución acumulada (CDF):} Comparación de CDFs observada y predicha, útil para verificar sesgo en percentiles.
    \item \textbf{Percentiles cruzados:} Tabla de P10, P25, P50, P75, P90, P95 observados vs predichos, con razones que diagnostican compresión o dilatación de rango.
\end{enumerate}

\section{Diagnósticos ejecutados}
\label{sec:diagnosticos}

Se desarrolló un programa de diagnósticos (\texttt{DIAGNÓSTICOS.py}) generando figuras y reportes estadísticos completos. Entre los diagnósticos clave se encuentran:

\begin{itemize}
    \item Dispersión de predicciones vs observaciones.
    \item Mapas espaciales de bias medio, RMSE, y correlación por píxel.
    \item Series temporales agregadas (promedio espacial) con intervalos intercuartílicos.
    \item Ciclo anual (promedio mensual climatológico).
    \item Métricas mensuales y anuales en tabla y gráficos.
    \item Análisis de extremos (percentiles altos).
    \item QQ-plot y función de distribución acumulada.
    \item Mapas de eficiencia (NSE, KGE) por píxel.
    \item Evolución temporal del error (RMSE y bias por paso temporal).
\end{itemize}

Los resultados se exportan tanto a figuras PNG (si se especifica \texttt{--save}) como a un archivo NetCDF con métricas espaciales (\texttt{*\_diagnosticos.nc}) para análisis posterior.

Cada una de las combinaciones de predictores descritas anteriormente fue evaluada múltiples veces mediante estas métricas, aleatorizando los meses de entrenamiento y validación, con el fin de concluir cuál era el conjunto de predictores más efectivo para el pronóstico. La elección de la alineación óptima se infirió con base en el script \texttt{COMPARACIÓN.py}.

\section{Procedimiento de predicción operativa}
\label{sec:prediccion_operativa}

Una vez entrenado y validado un modelo, se utiliza en modo predicción mediante el script \texttt{PRONÓSTICO.py}. El procedimiento es:

\begin{enumerate}
    \item \textbf{Cargar metadatos:} Se leen los parámetros de entrenamiento (normalización, PCA, transformación objetivo, etc.) desde \texttt{metadata.json} y los modelos PCA desde cache.
    \item \textbf{Validar predictores:} Se verifica que los archivos NetCDF nuevos posean la misma grilla, dimensionalidad y cobertura temporal que los datos de entrenamiento.
    \item \textbf{Procesar ventanas:} Para cada mes objetivo, se extrae una ventana de $L$ días previos, se normaliza, se aplica PCA, y se forma el tensor de entrada.
    \item \textbf{Predicción:} El modelo LSTM infiere sobre las secuencias procesadas, generando valores en el espacio transformado.
    \item \textbf{Inversión y clipping:} Las predicciones se transforman inversamente a milímetros, se clipean a rangos físicos plausibles, y se reconstruye la grilla 2D.
    \item \textbf{Exportación:} El resultado se guarda en formato NetCDF4 con metadatos descriptivos.
\end{enumerate}

\section{Procesamiento geográfico de resultados}
\label{sec:procesamiento_geografico}

Para análisis geográficos específicos (ej. desempeño en un departamento o región específica, mediante el paquete gtMapTools que incluye los shapes específicos del país utilizados en los análisis de INSIVUMEH), se implementó el script \texttt{DEP.py}, que:
\begin{enumerate}
    \item Carga un archivo de predicciones NetCDF (típicamente \texttt{modelo\_pixel\_predictions.nc}).
    \item Utiliza un archivo \textit{shape} de departamentos de Guatemala para definir polígonos.
    \item Construye una máscara espacial (píxel dentro/fuera del polígono) usando geometría computacional.
    \item Aplica la máscara a todos los pasos de tiempo de predicciones y observaciones.
    \item Exporta un NetCDF recortado conservando la dimensión temporal completa.
\end{enumerate}

\textbf{Este tipo de recorte geográfico es útil para validación regional y para productos operativos específicos de una zona. Por ejemplo, para propósitos de este documento, se consideró el \textit{shape} asociado a Bocacosta y se realizaron los diagnósticos de acuerdo a las variables de esa región. Y se utilizaron las métricas locales para elegir la combinación óptima de predictores para pronóstico de precipitación en la zona. 
}